% 附录 — Job Aid 速查卡（原书 PDF p.115–124）

\chapter{Job Aid 速查卡}

\label{app:job-aids}



\noindent\textit{Job Aids}



\begin{noteBox}

原书 PDF 第 115--124 页为折叠式速查卡，版式以图示与命令列表为主。

其中第 \textbf{115} 页（\emph{PrimeTime JOB AID: Timing Reports}）与第 \textbf{116} 页（\emph{PT JOB AID: Run Scripts}）

已从扫描页人工校对并译入下文；第 \textbf{121} 页（SI 串扰延迟微调）部分内容可辨认。

\textbf{第 117--120、122--123 页} OCR 质量极差，正文以图示速查为主——若需完整版式请对照原版扫描 PDF。

第 \textbf{124} 页为空白页（\textit{This page was intentionally left blank.}）。

\end{noteBox}



\section{PrimeTime JOB AID：Timing Reports}

\label{jobaid:timing-reports}



\noindent\textit{PrimeTime JOB AID: Timing Reports}（PDF p.115）



\subsection{实用 PrimeTime 命令}

\label{jobaid:timing-helpful}



\begin{description}[leftmargin=2.4em,style=nextline]

  \item[会话]%

\begin{lstlisting}

restore_session <unix_dir>   ;# 恢复前先 save_session

save_session <unix_dir>

\end{lstlisting}

  \item[帮助]%

\begin{lstlisting}

report_timing -help

help -verbose report_timing

man report_timing

help report*                 ;# 查找命令

apropos clock ;              ;# 搜索命令数据库

\end{lstlisting}

  \item[环境与历史]%

\begin{lstlisting}

printvar *sig*                ;# 查找变量

page_on                      ;# 别名

page_off                     ;# 别名

history

sh_list_key_bindings

set_app_var sh_line_editing_mode emacs

\end{lstlisting}

  \item[设计与库]%

\begin{lstlisting}

list_libraries ; list_designs

report_lib <lib_name>        ;# 时间单位

remove_design -all

remove_lib -all

\end{lstlisting}

\end{description}



\subsection{违例汇总报告}

\label{jobaid:timing-summary}



\begin{description}[leftmargin=2.8em,style=nextline]

  \item[\cmd{report\_analysis\_coverage}]%

  报告违例总数。

  \begin{itemize}

    \item \opt{-status\_violated}：违例详情

    \item \opt{-check \{setup hold\}}：指定检查类型

    \item \opt{-sort\_by slack}：默认排序

    \item \opt{-sort\_by check\_type}：按时钟组排序违例

  \end{itemize}

  \item[\cmd{report\_constraint -all}]%

  报告全部约束。

  \item[\cmd{report\_global\_timing}]%

  全局时序汇总。

  \item[\cmd{report\_qor -only\_violated}]%

  仅报告违例路径的 QoR。

\end{description}



\subsection{\cmd{report\_timing} 开关}

\label{jobaid:report-timing}



\subsubsection{控制路径数量与筛选}

\label{jobaid:rt-which}



\begin{itemize}

  \item \opt{-nworst}：每个端点考虑的路径条数

  \item \opt{-max\_paths}：每个设计生成的路径条数

  \item \opt{-delay}：setup（max）或 hold（min），及端点数据跳变\\

        \opt{max\_rise}、\opt{max\_fall}、\opt{min\_rise}、\opt{min\_fall}

  \item \opt{-group}：聚焦某 path group（如捕获时钟名）

  \item \opt{-group [get\_path\_group *]}

  \item 路径过滤：\opt{-to}、\opt{-rise\_to}、\opt{-fall\_to}、

        \opt{-from}、\opt{-rise\_from}、\opt{-fall\_from}、

        \opt{-through}、\opt{-rise\_through}、\opt{-fall\_through}、

        \opt{-exclude}、\opt{-rise\_exclude}、\opt{-fall\_exclude}、

        \opt{-slack\_lesser\_than}、\opt{-slack\_greater\_than}

\end{itemize}



\subsubsection{控制报告细节}

\label{jobaid:rt-details}



\begin{itemize}

  \item \opt{-input\_pins}：显示输入 pin，并分开网延迟与单元延迟

  \item \opt{-nets}：显示网名与 fanout

  \item \opt{-derate}：显示各延迟的 derate 因子

  \item \opt{-path\_full\_clock}：对 propagated clock 显示时钟网上的单元

  \item \opt{-path\_short}：从报告中去掉数据通路上的单元行

  \item \opt{-path\_end}：仅报告端点、数据到达、要求时间与 slack

  \item \opt{-crosstalk -nets -transition}：SI 延迟报告（含 delta delay）

  \item \opt{-pba\_mode path | exhaustive}：路径级重算（PBA）

\end{itemize}



\subsection{Tcl 示例}

\label{jobaid:rt-tcl}



\begin{itemize}

  \item \texttt{\#} 表示注释

  \item 命令名与开关可缩写（不必键入全名）

  \item 反斜杠 \texttt{\textbackslash} 用于续行

  \item 方括号 \texttt{[]} 用于嵌套命令

  \item 用 \cmd{get\_clocks} 引用时钟对象

\end{itemize}



\begin{lstlisting}

report_timing -fall_from [get_clocks PCI_CLK] \

  -delay min -max_paths 4

# 用 set_app_var 修改变量

set_app_var sh_enable_page_mode true

\end{lstlisting}



\section{PT JOB AID：Run Scripts}

\label{jobaid:run-scripts}



\noindent\textit{PT JOB AID: Run Scripts}（PDF p.116）



\subsection{实用 Tcl 过程（SolvNet）}

\label{jobaid:tcl-procs}



\begin{tabularx}{\textwidth}{@{}l X l@{}}

\toprule

\textbf{过程} & \textbf{说明} & \textbf{Doc ID} \\

\midrule

\cmd{aa} & always ask：同时搜索命令与变量 & 012959 \\

\cmd{view} & 在弹出窗口中查看长报告 & 014947 \\

\bottomrule

\end{tabularx}



\subsection{\cmd{.synopsys\_pt.setup} 示例}

\label{jobaid:pt-setup}



\begin{lstlisting}

# .synopsys_pt.setup

alias page_on {set_app_var sh_enable_page_mode true}

alias page_off {set_app_var sh_enable_page_mode false}

alias un_on {set_app_var timing_report_unconstrained_paths true}

alias un_off {set_app_var timing_report_unconstrained_paths false}

suppress_message {ENV-003 PTSR-004}

suppress_message CMD-029

history keep 200

set_app_var sh_enable_line_editing true

foreach _file [glob -nocomplain ./tcl_procs/*.tcl] {source $_file}

\end{lstlisting}



\subsection{\cmd{variables.tcl} 示例}

\label{jobaid:variables}



\begin{lstlisting}

# variables.tcl

set_app_var report_default_significant_digits 4

set_app_var link_create_black_boxes false

set_app_var sh_source_uses_search_path true

set_app_var sh_script_stop_severity E

set_app_var timing_update_status_level high

\end{lstlisting}



\subsection{查找非默认应用变量}

\label{jobaid:changed-vars}



\begin{lstlisting}

report_app_var -only_changed_vars

\end{lstlisting}



\subsection{Path-Based Analysis}

\label{jobaid:pba}



\begin{lstlisting}

report_timing -pba_mode path ...

report_timing -pba_mode exhaustive ...

\end{lstlisting}



\subsection{\cmd{RUN.tcl} 示例}

\label{jobaid:run-tcl}



\begin{lstlisting}

# RUN.tcl

start_profile

set_app_var search_path {. ./scripts ./libs ./designs}

lappend link_path mytech_lib.db RAM_lib.db

source -echo -verbose ./scripts/variables.tcl

read_verilog {top.v A.v B.v}

link_design TOP

read_parasitics ORCA.SPEF.gz

report_annotated_parasitics

source -echo -verbose constraints.tcl

redirect -tee ./EW.log {check_timing}

report_analysis_coverage

save_session top_savesession

redirect -tee -append ./EW.log {print_message_info}

redirect -tee -append ./EW.log {quit}

\end{lstlisting}



\subsection{从 Unix Shell 运行}

\label{jobaid:unix-run}



\begin{lstlisting}

unix% pt_shell -f RUN.tcl | tee -i run.log

# 若脚本异常终止，先修复错误

# 若脚本未正常结束，在 run.log 中搜索 Warning

# 可对消息编号生成 man 页以获取进一步信息

\end{lstlisting}



\section{其他 Job Aid 页（PDF p.117--124）}

\label{jobaid:other-pages}



\begin{noteBox}

\textbf{说明：}以下各页在原书中为图示速查卡；OCR 无法可靠还原全部版式。

条目根据 Lab 手册引用、可辨认扫描内容及 PrimeTime Workshop 惯例整理，供实验查阅。

对照原版 PDF 时请以扫描页为准。

\end{noteBox}



\subsection{Clocks and More（约 PDF p.117）}

\label{jobaid:clocks}



\noindent\textit{Job aid labeled ``Clocks and More''}（Lab~4 引用）



\begin{description}[leftmargin=2.6em,style=nextline]

  \item[\cmd{report\_clock}]%

  列出全部时钟（含 generated clock）、周期、波形、属性（propagated/ideal）及源 pin/port。

  \item[\cmd{report\_clock -attributes}]%

  显示时钟附加属性。

  \item[\cmd{report\_clock -skew [get\_clocks <clk>]}]%

  报告时钟不确定性及源延迟（setup/hold、early/late）。

  \item[\cmd{report\_clock\_timing}]%

  时钟网络时序（偏斜、延迟等）；SI 分析可加 \opt{-crosstalk}。

  \item[\cmd{get\_clocks} / \cmd{all\_clocks}]%

  获取时钟对象集合。

  \item[\cmd{check\_timing -override\_defaults clock\_crossing}]%

  报告时钟域交叉（与 GUI Clock Analyzer 矩阵对应）。

  \item[\cmd{set\_clock\_groups}]%

  定义时钟间关系（\opt{-asynchronous}、\opt{-logically\_exclusive}、\opt{-physically\_exclusive}）。

  \item[\cmd{create\_clock} / \cmd{create\_generated\_clock}]%

  创建主时钟与生成时钟。

  \item[\cmd{set\_propagated\_clock} / \cmd{set\_ideal\_network}]%

  切换 propagated 与 ideal 时钟网络。

\end{description}



\subsection{Constraints / Interface（约 PDF p.118--119）}

\label{jobaid:constraints}



\noindent\textit{约束与接口相关速查}（版式以原书图示为准）



\begin{itemize}

  \item \cmd{check\_timing -verbose}：约束完整性

  \item \cmd{report\_port -input\_delay} / \cmd{report\_port -output\_delay}：I/O 延迟约束

  \item \cmd{set\_input\_delay}、\cmd{set\_output\_delay}：端口延迟

  \item \cmd{set\_false\_path}、\cmd{set\_multicycle\_path}：例外路径

  \item \cmd{report\_constraint -all\_violators}：约束违例汇总

\end{itemize}



\subsection{PrimeTime SI：Crosstalk Delay Reports（约 PDF p.119--120）}

\label{jobaid:si-delay}



\noindent\textit{SI 延迟分析速查}（Lab~8 引用）



\begin{description}[leftmargin=2.8em,style=nextline]

  \item[启用 SI 延迟]%

\begin{lstlisting}

set si_enable_analysis true

read_parasitics -keep_coupling -format SBPF <file>

update_timing

\end{lstlisting}

  \item[\cmd{report\_timing -crosstalk}]%

  含 delta delay 的时序报告；常配合 \opt{-nets -transition -nosplit}。

  \item[\cmd{report\_si\_bottleneck}]%

  串扰瓶颈网；\opt{-cost\_type} 取值示例：

  \begin{itemize}

    \item \texttt{delta\_delay}：最大 delta delay（victim）

    \item \texttt{delta\_delay\_ratio}：delta/stage 比值（victim）

    \item \texttt{delay\_bump\_per\_aggressor}：相对 VDD 的 aggressor bump 之和

  \end{itemize}

  常用开关：\opt{-slack\_lesser\_than}、\opt{-include\_clock\_nets}、\opt{-max\_nets}。

  \item[\cmd{report\_clock\_timing -type skew -crosstalk -verbose}]%

  最坏时钟偏斜及串扰贡献。

  \item[修复思路]%

  \cmd{set\_coupling\_separation}（减小耦合）；加大 victim 驱动或插 buffer；缩小 aggressor 或隔离网线。

\end{description}



\subsection{PrimeTime SI Crosstalk Delay：Fine Tuning（PDF p.121）}

\label{jobaid:si-finetune}



\noindent\textit{PrimeTime SI Crosstalk Delay: Fine Tuning for Timing Closure}



\paragraph{SolvNet 文章}

\begin{tabularx}{\textwidth}{@{}X r@{}}

\toprule

\textbf{主题} & \textbf{Doc ID} \\

\midrule

在网上生成延时计算报告的简便方法 & 013383 \\

有效 aggressor 如何因 bump 大小被筛除 & 014179 \\

PrimeTime path-based analysis 精确 sign-off & 012134 \\

在 shell 中用弹出窗口查看长报告（\cmd{view}） & 014947 \\

用 Tcl 过程 \cmd{aa} 搜索 Synopsys 命令与变量 & 012959 \\

\bottomrule

\end{tabularx}



\paragraph{控制特定网的耦合信息}

\begin{lstlisting}

set_si_delay_analysis

\end{lstlisting}



\paragraph{What-If / ECO 相关命令}

\begin{lstlisting}

size_cell

report_lib tech_lib buf*

insert_buffer

remove_buffer

set_coupling_separation

set_eco_instance_name_prefix

set_eco_net_name_prefix

estimate_eco

fix_eco_timing

fix_eco_drc

write_changes

\end{lstlisting}



\paragraph{PBA 相关命令}

\begin{lstlisting}

set viol_path [get_timing_paths -slack_lesser_than 0]

set path [get_timing_paths -pba_mode exhaustive $viol_path]

report_timing -crosstalk $path

report_timing -pba_mode path ...

report_timing -pba_mode exhaustive ...

\end{lstlisting}



\paragraph{调试变量}

\begin{lstlisting}

# Variables to control victim reselection

printvar si_xtalk_reselect*

# Variables to control exit criteria

printvar si_xtalk_exit*

# Variable to control electrical filtering

printvar si_filter*

\end{lstlisting}



\subsection{PrimeTime SI：Noise Analysis（约 PDF p.122--123）}

\label{jobaid:si-noise}



\noindent\textit{SI 噪声分析速查}（Lab~9 引用）



\begin{description}[leftmargin=2.8em,style=nextline]

  \item[流程]%

  \cmd{si\_enable\_analysis true} $\rightarrow$ \cmd{check\_noise} $\rightarrow$ \cmd{update\_noise} $\rightarrow$ \cmd{report\_noise}

  \item[\cmd{set\_noise\_parameters}]%

  默认开启 between-the-rails 分析与 timing window 滤波 aggressor。

  \item[\cmd{report\_noise}]%

  \opt{-all\_violators}、\opt{-slack\_type area | area\_percent}、\opt{-data\_pins}、\opt{-clock\_pins}、\opt{-nworst}。

  \item[\cmd{report\_noise\_calculation}]%

  单网详细噪声计算（\opt{-above\_low} / \opt{-below\_high}、\opt{-from}、\opt{-to}）。

  \item[\cmd{set\_si\_noise\_analysis}]%

  排除 aggressor 或定义无限 timing window。

  \item[\cmd{set\_noise\_margin}]%

  用户自定义噪声裕量。

  \item[\cmd{report\_noise\_bumps}]%

  按 bump 高度筛选（如 \opt{-bump\_greater\_than 0.5}）。

\end{description}



\subsection{空白页（PDF p.124）}

\label{jobaid:blank}



原书最后一页为刻意留白（\textit{This page was intentionally left blank.}）。


